\section{微积分中的一些结果}
\subsection{微分}
\textbf{定义}$^\clubsuit $:令$A\subset\mathbb{R}^n$为开集,$f:A\to \mathbb{R}^m$,设$A$包含$a$的一个邻域.如果有线性映射$L:\mathbb{R}^n\to \mathbb{R}^m$使得
$$
\lim\limits_{h\to 0}\frac{|f(a+h)-f(a)-L(h)|}{|h|}=0
$$
则称$f$在$a$处可微,并且称$L$为$f$在$a$处的微分,记为$Df(a)$.如果记$f=(f_1,\cdots,f_m)$,那么$Df(a)$是一个$m\times n$矩阵,即
$$
Df(a)=
\begin{pmatrix}
    \frac{\partial f_1}{\partial x_1}(a) & \cdots & \frac{\partial f_1}{\partial x_n}(a)\\
    \vdots & \ddots & \vdots\\
    \frac{\partial f_m}{\partial x_1}(a) & \cdots & \frac{\partial f_m}{\partial x_n}(a)\\
\end{pmatrix}
$$
\begin{tcolorbox}[colback=KuangNei,colframe=BianKuang,title=\textbf{微分中值定理},sharpish corners]
    设 $D\subset\mathbb{R}^{n}$ 为凸域，函数 $f:D\rightarrow\mathbb{R}$在$D$中处处可微,则任给 $x,y\in D$,存在$\theta\in(0,1)$,使得  
$$
f(x)-f(y)=\nabla f(\xi)\cdot(x-y),\quad\xi=\theta x+(1-\theta)y.
$$
\end{tcolorbox}
\textbf{Proof}:令$\sigma(t)=t x+(1-t)y$ .由$D$为凸域可知当 $t\in[0,1]$ 时 $\sigma(t)\in D$.对一元函数$\varphi(t)=f(\sigma(t))$用微分中值定理可知,存在 $\theta\in(0,1)$ ,使得 $\varphi(1)-\varphi(0)=\varphi^{\prime}(\theta)$.由复合函数求导可得  
$$
\varphi(1)-\varphi(0)=\nabla f(\xi)\cdot\sigma^{\prime}(\theta)=\nabla f(\xi)\cdot(x-y),
$$
其中$\xi=\sigma(\theta)=\theta x+(1-\theta)y.$由$f(x)=\varphi(1)$,$f(y)=\varphi(0)$ 可知欲证结论成立. 
\begin{tcolorbox}[colback=KuangNei,colframe=BianKuang,title=\textbf{拟微分中值定理},sharpish corners]
设 $D\subset\mathbb{R}^{n}$为凸域,$f:D\rightarrow\mathbb{R}^{m}$ 在 $D$ 中处处可微,则任给 $x,y\in D$ ,存在 $\xi\in D$ ,使得 
$$
\|f(x)-f(y)\|\leqslant\|Df(\xi)\|\cdot\|x-y\|.
$$
\end{tcolorbox}
\textbf{Proof}:基本的想法是对 $f$ 的分量的线性组合应用微分中值定理,为此,不妨设 $f\left(x\right)\neq f\left(y\right)$ 任意取定 $\mathbb{R}^{m}$ 中的单位向量 $u=(u_{1},u_{2},\cdots,u_{m})$ ,记  
$$
g=u\cdot f=\sum_{i=1}^{m}u_{i}f_{i},
$$
则 $g$为$D$中可微函数.根据微分中值定理,存在 $\xi\in D$ ,使得  
$$
g(x)-g(y)=\nabla g(\xi)\cdot(x-y).
$$
注意到 $\nabla g(\xi)=\sum_{i=1}^{m}u_{i}\nabla f_{i}(\xi)$ ,利用Cauchy-Schwarz不等式可得  
$$
\|\nabla g(\xi)\|\leqslant\sum_{i=1}^{m}|u_{i}|\cdot\|\nabla f_{i}(\xi)\|\leqslant\|u\|\cdot\Big(\sum_{i=1}^{m}\|\nabla f_{i}(\xi)\|^{2}\Big)^{\frac{1}{2}}=\|Df(\xi)\|.
$$
由 $g(x)-g(y)=u\cdot\left[f(x)-f(y)\right]$ 可得 
$$
\left|u\cdot[f(x)-f(y)]\right|\leqslant\|\nabla g(\xi)\|\cdot\|x-y\|\leqslant\|Df(\xi)\|\cdot\|x-y\|.
$$
在上式中取 $u={\frac{f(x)-f(y)}{\|f(x)-f(y)\|}}$ 就完成了定理的证明.
\begin{tcolorbox}[colback=KuangNei,colframe=BianKuang,title=\textbf{逆映射定理},sharpish corners]
    设 $D\subset\mathbb{R}^{n}$ 为开集, $f:D\to\mathbb{R}^{n}$ 为 $C^{k}(k\geqslant1)$ 映射,$x^{0}\in D$ 如果 $\operatorname*{det}J f(x^{0})\neq0$ ,则存在 $x^{0}$ 的开邻域 $U\subset D$ 
    以及 $y^{0}=f(x^{0})$ 的开邻域 $V\subset\mathbb{R}^{n}$ 使得 $f|_{U}:U\to V$ 是可逆映射,且其逆仍为 $C^{k}$ 映射.
\end{tcolorbox}
\textbf{Proof}:不失一般性,可设 $x^{0}=0$ $y^{0}=0$ 以 $L$ 记 $f$ 在 $x^{0}=0$ 处的微分,则 $L$ 可逆,且 $L^{-1}\circ f$ 在 $x^{0}$ 处的微分为恒同映射.如果欲证结论对 $L^{-1}\circ f$ 成立,则对 $f$ 也成立.因此,不妨从一开始就假设 $D f(x^{0})=I_{n}$  
在 $x^{0}=0$ 附近, $f$ 是恒同映射的小扰动:  
$$
f(x)=x+g(x),D g(0)=0.
$$
扰动项 $g(x)=f(x)-x$ 为 $C^{k}$ 映射,因此,存在 $\delta>0$ 使得  
$$
\|D g(x)\|\leqslant\frac{1}{2},~\forall~x\in\overline{{B_{\delta}(0)}}\subset D.
$$
由拟微分中值定理可知  
$$
\|g(x_{1})-g(x_{2})\|\leqslant\frac{1}{2}\|x_{1}-x_{2}\|,\forall x_{1},x_{2}\in\overline{{B_{\delta}(0)}}.
$$
给定 $y\in B_{\frac{\delta}{2}}(0)$ ,我们在 $B_{\delta}(0)$ 中解方程  
\begin{equation}
    f(x)=y,\mathrm{\mathrm{~}}\mathrm{\mathrm{H~}}x=y-g(x).
    \label{eq:11.16}
\end{equation}
记 $\varphi(x)=y-g(x)$ 当 $x\in\overline{{B_{\delta}(0)}}$ 时，  
\begin{equation}
    \|\varphi(x)\|\leqslant\|y\|+\|g(x)\|<\frac{\delta}{2}+\frac{1}{2}\|x\|\leqslant\delta.
    \label{eq:11.17}
\end{equation}
这说明 $\varphi\big(\overline{{B_{\delta}(0)}}\big)\subset B_{\delta}(0)$ 当 $x_{1},x_{2}\in\overline{{B_{\delta}(0)}}$ 时,  
$$
\|\varphi(x_{1})-\varphi(x_{2})\|=\|g(x_{2})-g(x_{1})\|\leqslant{\frac{1}{2}}\|x_{1}-x_{2}\|.
$$
根据压缩映射原理,方程\eqref{eq:11.16}在 $\overline{{B_{\delta}(0)}}$ 中有唯一解,记为$x_{y}$:由\eqref{eq:11.17}式可知 $x_{y}\in B_{\delta}(0)$ 令  
$$
U=f^{-1}\Big(B_{\frac{\delta}{2}}(0)\Big)\cap B_{\delta}(0),V=B_{\frac{\delta}{2}}(0),
$$
则我们已经证明了 $f|_{U}:U\to V$ 是一一映射,其逆映射 $h(y)=x_{y}$ 满足  
$$
y-g(h(y))=h(y).
$$
(1) $h:V\to U$ 是连续映射:当 $y_{1}$ $y_{2}\in V$ 时,  
$$
h(y_{1})-h(y_{2})\|\leqslant\|y_{1}-y_{2}\|+\|g(h(y_{1}))-g(h(y_{2}))\|\leqslant\|y_{1}-y_{2}\|+\frac{1}{2}\|h(y_{1})-h(y_{2})\|
$$
这说明 $\|h(y_{1})-h(y_{2})\|\leqslant2\|y_{1}-y_{2}\|,\forall y_{1},y_{2}\in V.$\\  
(2) $h:V\to U$ 是可微映射:设 $y_{0}\in V$ ,则对 $y\in V$ 有  
$$
\begin{aligned}
    h(y)-h(y_{0})&=(y-y_{0})-\big[g(h(y))-g(h(y_{0}))\big]\\
                 &=(y-y_{0})-D g(h(y_{0}))\cdot(h(y)-h(y_{0}))+o\big(\|h(y)-h(y_{0})\|\big).
\end{aligned}
$$
利用$D f=I_{n}+D g$和(1),上式可改写为  
$$
D f(h(y_{0}))\cdot(h(y)-h(y_{0}))=(y-y_{0})+o\big(\|y-y_{0})\|\big),
$$
因而  
$$
h(y)-h(y_{0})=\left[D f(h(y_{0}))\right]^{-1}\cdot(y-y_{0})+o\big(\|y-y_{0}\|\big).
$$
这说明 $h$ 在 $y_{0}$ 处可微.\\  
(3) $h:V\to U$ 为 $C^{k}$ 映射:由(2)可知  
\begin{equation}
    D h(y)=\left[D f(h(y))\right]^{-1},~\forall y\in V.
    \label{eq:11.19}
\end{equation}
由 $f\in C^{k}$ 知 $D f\in C^{k-1}$.由(2)及\eqref{eq:11.19}式可推出 $D h$ 连续,即 $h\in C^{1}$ ,再由 $D f\in C^{k-1}$ $h\in C^{1}$ 及\eqref{eq:11.19}式可推出 $D h\in C^{1}$ 即 $h\in C^{2}$ .依此类推,最后我们就得到 $h\in C^{k}$.\par
\textbf{命题}$^\clubsuit $:设函数$f$在凸域$D$中处处可微,那么\\  
(i)$f$为凸函数$\Longleftrightarrow$  
\begin{equation}
f(y)\geqslant f(x)+\nabla f(x)\cdot(y-x),\forall x,y\in D.
\label{eq:11.20}
\end{equation}
(ii)当$f\in C^{2}(D)$时,$f$为凸函数$\Longleftrightarrow\nabla^{2}f$(半正定).\par
\textbf{Proof}: (i)$"\Longrightarrow"$任给$x,y\in D,t\in(0,1)$ ,有  
$$
t\big[f(y)-f(x)\big]\geqslant f(x+t(y-x))-f(x)=\nabla f(x)\cdot t(y-x)+o\big(t\|y-x\|\big),
$$
上式两边除以 $t$ ，然后令 $t\to0^{+}$ 即得\eqref{eq:11.20}式  
$\Longleftarrow^{\mathfrak{N}}$ 任给 $x,y\in D$ $t\in(0,1)$ ，记 $z=t x+(1-t)y,$ 则  
$$
f(x)\geqslant f(z)+\nabla f(z)\cdot(x-z),f(y)\geqslant f(z)+\nabla f(z)\cdot(y-z).
$$
这说明  
$$
t f(x)+(1-t)f(y)\geqslant f(z)+\nabla f(z)\cdot\left[t(x-z)+(1-t)(y-z)\right]=f(z).
$$
(ii) $"\Longleftarrow"$设 $f\in C^{2}(D)$ ,根据之前的讨论,任给 $x,y\in D$ ,存在 $\xi\in D_{\mathrm{i}}$ 使得  
$$
f(\boldsymbol{y})=f(\boldsymbol{x})+\nabla f(\boldsymbol{x})\cdot(\boldsymbol{y}-\boldsymbol{x})+\frac{1}{2}(\boldsymbol{y}-\boldsymbol{x})\nabla^{2}f(\boldsymbol{\xi})(\boldsymbol{y}-\boldsymbol{x})^{\mathrm{T}}.
$$
由 $\nabla^{2}f\geqslant0$ 知\eqref{eq:11.20}式成立,再由(i)知 $f$ 为凸函数\par  
$"\Longrightarrow"$(反证法):如果$\nabla^{2}f(x)$不是半正定的,则存在 $\mathbb{R}^{n}$ 中的单位向量 $h$, 使得 $h\nabla^{2}f(x)h^{\mathrm{T}}$ $<0$ 对一元函数 $\varphi(t)=f(x+t h)$ 应用Taylor公式可得  
$$
f(x+t h)=f(x)+t\nabla f(x)\cdot h+{\frac{1}{2}}t^{2}h\nabla^{2}f(x)h^{\mathrm{T}}+o(t^{2})\quad(t\to0).
$$
当 $t\neq0$ 充分小时,上式右端第三项小于零,这与\eqref{eq:11.20}式相矛盾.\par
\begin{tcolorbox}[colback=KuangNei,colframe=BianKuang,title=\textbf{Taylor公式},sharpish corners]
设 $D\subset\mathbb{R}^{n}$为凸域, $f\in C^{m+1}(D)$ ,$a\in D$ ,则任给 $x\in D,$ 存在 $\theta\in(0,1)$ 使得  
$$
f(x)=\sum_{|\alpha|\leqslant m}{\frac{\mathrm{D}^{\alpha}f(a)}{\alpha!}}(x-a)^{\alpha}+\sum_{|\alpha|=m+1}{\frac{\mathrm{D}^{\alpha}f(a+\theta(x-a))}{\alpha!}}(x-a)^{\alpha}.
$$
\end{tcolorbox}
\textbf{例}$^\spadesuit$ :设$f:\mathbb{R}^2\to\mathbb{R}$为$ C^k(k\geqslant 1)$函数,$\frac{\partial f}{\partial x}(x^0,y^0)\neq 0$,在$(x^0,y^0)$附近解方程 $f(x,y)=f(x_{0},y_{0})$.\par
显然,$(x_{0},y_{0})$ 是方程的解,利用反函数定理,我们可以在 $(x_{0},y_{0})$ 附近找到别的解.为此,令  
$$
F:\mathbb{R}^{2}\to\mathbb{R}^{2},F(x,y)=(x,f(x,y)).
$$
在 $(x_{0},y_{0})$ 处，有  
$$
\operatorname*{det}D F(x^{0},y^{0})=\left|{\begin{array}{c c}{1}&{0}\\ {{\cfrac{\partial f}{\partial x}}(x^{0},y^{0})}&{{\cfrac{\partial f}{\partial y}}(x^{0},y^{0})}\end{array}}\right|={\cfrac{\partial f}{\partial y}}(x^{0},y^{0})\neq0.
$$
由反函数定理,在 $(x^{0},y^{0})$ 附近 $F$ 为可逆映射.于是当 $x$ 在 $x^{0}$ 附近时,记 
$$
F^{-1}\bigl(x,f(x^{0},y^{0})\bigr)=\bigl(\varphi(x),\psi(x)\bigr),
$$
$\varphi(x),\psi(x)$ 均为 $C^{k}$ 函数.根据 $F$ 的定义可得  
$$
\bigl(\varphi(x),f(\varphi(x),\psi(x))\bigr)=\bigl(x,f(x^{0},y^{0})\bigr),
$$
这说明 $\varphi(x)=x,f(x,\psi(x))=f(x^{0},y^{0})$ 对 $x$ 求导还可得到  
$$
{\frac{\partial f}{\partial x}}(x,\psi(x))+{\frac{\partial f}{\partial y}}\big(x,\psi(x)\big)\cdot\psi^{\prime}(x)=0,
$$
从而有  
$$
\psi^{\prime}(x)=-\Big[\frac{\partial f}{\partial y}(x,\psi(x))\Big]^{-1}\frac{\partial f}{\partial x}(x,\psi(x)).
$$
$y=\psi(x)$ 称为由 $f(x,y)=f(x_{0},y_{0})$ 所决定的隐函数.
\begin{tcolorbox}[colback=KuangNei,colframe=BianKuang,title=\textbf{隐映射定理},sharpish corners]
设 $W\subset\mathbb{R}^{n+m}$ 为开集, $W$ 中的点用 $(x,y)$ 表示,其中 $x=$ $\left(x_{1},x_{2},\cdots,x_{n}\right)$ ,$y=(y_{1},y_{2},\cdots,y_{m}).f:W\to\mathbb{R}^{m}$ 为 $C^{k}$ 映射,用分量表示为  
$$
f(x,y)=(f_{1}(x,y),f_{2}(x,y),\cdots,f_{m}(x,y)).
$$
设 $(x^{0},y^{0})\in W$ ,且$\operatorname*{det}D_{y}f(x^{0},y^{0})\neq0$ ,其中 $\begin{array}{r}{D_{y}f(x,y)=\left(\frac{\partial f_{i}}{\partial y_{j}}(x,y)\right)_{m\times m}}\end{array}$ 则存在 $x^{0}$ 的开邻域 $V\subset\mathbb{R}^{n}$ 以及唯一的 $C^{k}$ 映射 $\psi:V\to{\mathbb{R}}^{m}$ 使得\par  
(1)$y^{0}=\psi(x^{0}),f(x,\psi(x))=f(x_{0},y_{0}),\forall x\in V;$\par
(2)$J\psi(x)=-\big[J_{y}f(x,\psi(x))\big]^{-1}D_{x}f(x,\psi(x))$,其中$D_{x}f(x,y)=\Big(\frac{\partial f_{i}}{\partial x_{j}}(x,y)\Big)_{m\times n}.$
\end{tcolorbox}
\subsection{积分}
\textbf{散度定理}$^\clubsuit $:设$\Omega\subset\mathbb{R}^n$是一个具有$C^1$边界的有界区域,并设$\mathbf{\nu} $表示$\partial\Omega$的单位外法向量.于是对$C^1(\overline{\Omega})$中的任一向量场$\mathbf{F}$有
$$
\int_{\Omega}\nabla\cdot\mathbf{F}dx=\int_{\partial\Omega}\mathbf{F}\cdot\mathbf{\nu}dS.
$$
如果$u$是$C^2(\overline{\Omega})$函数,那么取$\mathbf{F}=Du$,则有
$$
\int_{\Omega}\Delta u dx=\int_{\partial\Omega}Du\cdot\mathbf{\nu}dS=\int_{\partial\Omega}\frac{\partial u}{\partial\mathbf{\nu}}dS.
$$
由散度定理,很容易得到所谓的分部积分公式:设$u,v\in C^1(\overline{\Omega})$,则
$$
\int_{\Omega}vDu\mathrm{d}x=-\int_{\Omega}uDv\mathrm{d}x+\int_{\partial\Omega}uv\mathbf{\nu}\mathrm{d}S.
$$
\begin{tcolorbox}[colback=KuangNei,colframe=BianKuang,title=\textbf{Green恒等式},sharpish corners]
    设$u,v\in C^2(\overline{\Omega})$,则\\
    1.Green第一恒等式
    $$
    \int_{\Omega}v\Delta u dx+\int_{\Omega}Du\cdot Dvdx=\int_{\partial\Omega}v\frac{\partial u}{\partial \mathbf{\nu}}dS.
    $$
    2.Greend第二恒等式
    $$
    \int_{\Omega}\left(v\Delta u-u\Delta v\right)dx=\int_{\partial\Omega}\left(v\frac{\partial u}{\partial \mathbf{\nu}}-u\frac{\partial v}{\partial \mathbf{\nu}}\right)dS.
    $$
\end{tcolorbox}






